\chapter{Analyse des besoins et conception}
\label{chap:conception}

Ce chapitre précise les besoins de Tastify et présente la conception fonctionnelle et objet. La modélisation s'appuie sur UML, langage standard pour représenter les systèmes logiciels \cite{uml}. Les diagrammes présentés décrivent les cas d'utilisation, la structure de données et les interactions dynamiques du projet.

\section{Acteurs du système}
Tastify distingue quatre acteurs fonctionnels. Les accès dépendent du rôle associé à chaque utilisateur.

\begin{table}[H]
\centering
\caption{Acteurs et responsabilités}
\label{tab:acteurs}
\begin{tabularx}{\textwidth}{p{3cm}Y}
\toprule
\textbf{Acteur} & \textbf{Responsabilités principales} \\
\midrule
Gérant & Gérer le restaurant, gérer le menu, les stocks, les employés, les paramètres, les réservations, les avis et les indicateurs. \\
Serveur & Consulter le plan de salle, créer les commandes, suivre les tables et déclencher l'encaissement. \\
Cuisinier & Consulter les commandes cuisine, mettre à jour les statuts de préparation et signaler les plats prêts. \\
Client & Consulter le menu, s'inscrire, se connecter, réserver, payer, consulter sa fidélité et laisser un avis. \\
\bottomrule
\end{tabularx}
\end{table}

\section{Besoins fonctionnels}
Les besoins fonctionnels sont regroupés par domaine métier dans le tableau~\ref{tab:besoins-fonctionnels}.

\begin{longtable}{p{3.2cm}p{10.5cm}}
\caption{Besoins fonctionnels principaux}
\label{tab:besoins-fonctionnels}\\
\toprule
\textbf{Module} & \textbf{Besoins} \\
\midrule
\endfirsthead
\toprule
\textbf{Module} & \textbf{Besoins} \\
\midrule
\endhead
Authentification & Inscription et connexion, séparation des portails staff/client, jetons JWT, refresh par cookie, contrôle des accès par rôle. \\
Menu & Créer, modifier, masquer ou afficher des catégories et des plats ; gérer la disponibilité et les images. \\
Salle & Visualiser les tables, leur capacité et leur statut ; accéder à la prise de commande par table. \\
Commandes & Créer des commandes, enregistrer les lignes, figer les prix unitaires et suivre les statuts de préparation. \\
Cuisine & Afficher les commandes dans un KDS et synchroniser les changements de statut vers le staff. \\
Stock & Gérer les ingrédients, les seuils d'alerte, les recettes et les mouvements de stock. \\
Réservations & Permettre au client de réserver et au staff de suivre les demandes selon la disponibilité. \\
Paiements & Enregistrer les paiements, gérer les méthodes disponibles et associer des contributions aux lignes de commande. \\
Fidélité & Suivre les points, transactions, paliers et récompenses clients. \\
Avis & Collecter les commentaires, stocker une note optionnelle et déclencher l'analyse de sentiment. \\
Analytics & Agréger des indicateurs : revenus, commandes, plats populaires, stocks et sentiment client. \\
Configuration & Paramétrer les informations du restaurant, la devise, certains réglages de service et l'identité visuelle. \\
\bottomrule
\end{longtable}

\section{Besoins non fonctionnels}
Tastify doit aussi respecter plusieurs contraintes techniques pour rester maintenable.

\begin{table}[H]
\centering
\caption{Besoins non fonctionnels}
\label{tab:besoins-non-fonctionnels}
\begin{tabularx}{\textwidth}{p{3.2cm}Y}
\toprule
\textbf{Besoin} & \textbf{Traduction dans le projet} \\
\midrule
Sécurité & Authentification JWT, permissions DRF, rôles gérant/serveur/cuisinier/client, séparation des portails. \\
Maintenabilité & Découpage en applications Django spécialisées et deux frontends React séparés. \\
Réactivité & WebSocket pour les événements staff et KDS ; tâches Celery pour les traitements non bloquants. \\
Traçabilité & Modèles avec dates de création et de mise à jour ; conservation des prix de commande. \\
Déploiement & Docker Compose avec services backend, base de données, Redis, worker, beat et frontends. \\
Testabilité & Tests backend avec pytest, tests frontend avec Vitest et scénarios end-to-end avec Playwright. \\
\bottomrule
\end{tabularx}
\end{table}

\section{Diagrammes de cas d'utilisation}
Les diagrammes de cas d'utilisation décrivent les échanges principaux entre les acteurs et les modules.

\begin{figure}[H]
\centering
\includegraphics[width=\textwidth]{figure_08.jpeg}
\caption{Diagramme de cas d'utilisation du back-office gérant}
\label{fig:uc-admin}
\end{figure}

La figure~\ref{fig:uc-admin} regroupe les fonctions d'administration : menu, stock, employés, alertes, tableaux de bord et rapports.

\begin{figure}[H]
\centering
\includegraphics[width=\textwidth]{figure_05.png}
\caption{Diagramme de cas d'utilisation des modules salle et cuisine}
\label{fig:uc-salle-cuisine}
\end{figure}

La figure~\ref{fig:uc-salle-cuisine} montre le lien entre le serveur, la salle et le KDS. Elle explique le besoin de temps réel : la commande doit passer de l'interface serveur vers l'écran cuisine sans attente inutile.

\begin{figure}[H]
\centering
\includegraphics[width=0.98\textwidth]{figure_02.png}
\caption{Diagramme de cas d'utilisation du portail client}
\label{fig:uc-client}
\end{figure}

La figure~\ref{fig:uc-client} présente les fonctions disponibles côté client : consultation du menu, réservation, commande, paiement, fidélité et avis.

\section{Diagramme de classes}
Le diagramme de classes présente les entités majeures et leurs relations.

\clearpage
\begin{landscape}
\begin{figure}[p]
\centering
\includegraphics[width=0.95\linewidth,height=0.82\textheight,keepaspectratio]{figure_03.png}
\caption{Diagramme de classes de Tastify}
\label{fig:classes}
\end{figure}
\end{landscape}
\clearpage

La figure~\ref{fig:classes} correspond aux domaines présents dans l'application : utilisateurs, plats, tables, commandes, réservations, paiements, avis et stock. Certaines entités détaillées dans le code ne sont pas toutes visibles sur le schéma structurel initial ; elles sont donc précisées dans le modèle logique ci-dessous.

\section{Diagrammes de séquence}
Afin d'illustrer la dynamique des échanges et le comportement du système lors de scénarios clés, trois diagrammes de séquence sont introduits ci-dessous.

\subsection{Authentification staff}
Le processus d'authentification sécurisé pour les membres du personnel repose sur des jetons JWT. Le serveur, le cuisinier ou le gérant fournit ses identifiants qui sont validés par l'API backend avant d'émettre les clés de session, comme le montre la figure~\ref{fig:seq-auth}.

\begin{figure}[H]
\centering
\resizebox{\textwidth}{!}{%
\begin{tikzpicture}[>=stealth]
  \colorlet{pboxfill}{amber!10}
  \colorlet{pboxdraw}{amber!80!black}
  \colorlet{actfill}{gray!20}
  \colorlet{actdraw}{gray!60}
  \colorlet{msgcolor}{blue!60!black}
  \tikzset{
    participant/.style={draw=pboxdraw, fill=pboxfill, rounded corners=2pt, minimum width=2.8cm, minimum height=0.65cm, align=center, font=\bfseries\small},
    activation/.style={draw=actdraw, fill=actfill},
    msg/.style={thick, msgcolor},
    ret/.style={thick, dashed, gray!80!black},
    note/.style={midway, above, fill=white, inner sep=1.5pt, font=\scriptsize}
  }

  \foreach \x/\name in {0/Personnel,4.3/Interface Staff,8.6/API Backend,12.9/Base données} {
    \node[participant] at (\x,0) {\name};
    \draw[dashed, thick, gray!55] (\x,-0.45) -- (\x,-6.35);
  }

  \draw[activation] (-0.09,-0.8) rectangle (0.09,-1.25);
  \draw[activation] (4.21,-1.05) rectangle (4.39,-2.55);
  \draw[activation] (8.51,-2.25) rectangle (8.69,-5.2);
  \draw[activation] (12.81,-3.25) rectangle (12.99,-4.25);
  \draw[activation] (4.21,-4.75) rectangle (4.39,-5.8);

  \draw[->, msg] (0.12,-1.05) -- node[note, align=center] {1. Identifiants\\nom + code} (4.18,-1.05);
  \draw[->, msg] (4.42,-2.05) -- node[note] {2. POST \code{/api/token/}} (8.48,-2.05);
  \draw[->, msg] (8.72,-3.05) -- node[note] {3. Vérifier rôle} (12.78,-3.05);
  \draw[<-, ret] (8.72,-4.05) -- node[note] {4. Profil validé} (12.78,-4.05);
  \draw[<-, ret] (4.42,-5.0) -- node[note, align=center] {5. Token\\+ cookie} (8.48,-5.0);
  \draw[<-, msg] (0.12,-5.65) -- node[note] {6. Accès dashboard} (4.18,-5.65);
\end{tikzpicture}
}
\caption{Diagramme de séquence de l'authentification staff}
\label{fig:seq-auth}
\end{figure}

La figure~\ref{fig:seq-auth} détaille la séparation entre la saisie côté interface, la validation par l'API et la vérification du rôle en base de données avant l'accès au tableau de bord.

\subsection{Prise de commande et envoi vers le KDS}
Ce scénario décrit le flux de traitement en temps réel depuis la prise de commande par le serveur en salle jusqu'à sa diffusion instantanée vers l'écran de cuisine KDS grâce au protocole WebSocket géré par Redis, comme le montre la figure~\ref{fig:seq-commande-kds}.

\begin{figure}[H]
\centering
\resizebox{\textwidth}{!}{%
\begin{tikzpicture}[>=stealth]
  \colorlet{pboxfill}{amber!10}
  \colorlet{pboxdraw}{amber!80!black}
  \colorlet{actfill}{gray!20}
  \colorlet{actdraw}{gray!60}
  \colorlet{msgcolor}{blue!60!black}
  \tikzset{
    participant/.style={draw=pboxdraw, fill=pboxfill, rounded corners=2pt, minimum width=2.65cm, minimum height=0.65cm, align=center, font=\bfseries\small},
    activation/.style={draw=actdraw, fill=actfill},
    msg/.style={thick, msgcolor},
    ret/.style={thick, dashed, gray!80!black},
    note/.style={midway, above, fill=white, inner sep=1.5pt, font=\scriptsize}
  }

  \foreach \x/\name in {0/Serveur,3.65/Interface,7.3/API Backend,10.95/Serveur WS,14.6/KDS Cuisine} {
    \node[participant] at (\x,0) {\name};
    \draw[dashed, thick, gray!55] (\x,-0.45) -- (\x,-7.25);
  }

  \draw[activation] (-0.09,-0.75) rectangle (0.09,-1.2);
  \draw[activation] (3.56,-1.0) rectangle (3.74,-4.2);
  \draw[activation] (7.21,-1.95) rectangle (7.39,-5.2);
  \draw[activation] (10.86,-4.85) rectangle (11.04,-6.2);
  \draw[activation] (14.51,-5.85) rectangle (14.69,-7.0);

  \draw[->, msg] (0.12,-1.0) -- node[note] {1. Table et plats} (3.53,-1.0);
  \draw[->, msg] (3.77,-1.95) -- node[note] {2. Créer commande} (7.18,-1.95);
  \draw[->, msg] (7.42,-2.85) -- ++(0.75,0) -- ++(0,-0.55) -- node[right, align=left, fill=white, inner sep=1.5pt, font=\scriptsize, pos=0.55] {3. Sauvegarde\\en base} (7.42,-3.4);
  \draw[<-, ret] (3.77,-4.0) -- node[note] {4. Commande créée} (7.18,-4.0);
  \draw[->, msg] (7.42,-4.95) -- node[note] {5. Notification staff} (10.83,-4.95);
  \draw[->, msg] (11.07,-5.95) -- node[note] {6. Broadcast WS} (14.48,-5.95);
  \draw[->, msg] (14.6,-6.75) -- ++(0.75,0) -- ++(0,-0.35) -- node[right, align=left, fill=white, inner sep=1.5pt, font=\scriptsize, pos=0.45] {7. Affichage\\temps réel} (14.6,-7.1);
\end{tikzpicture}
}
\caption{Diagramme de séquence de la prise de commande et envoi au KDS}
\label{fig:seq-commande-kds}
\end{figure}

La figure~\ref{fig:seq-commande-kds} met en évidence le passage de la commande depuis l'interface staff vers l'API, puis la notification WebSocket qui permet au KDS cuisine de se mettre à jour en temps réel.

\subsection{Création d'un avis client et analyse de sentiment}
Le client soumet son commentaire et sa note depuis le portail public. Le backend enregistre l'avis et délègue en tâche d'arrière-plan Celery l'appel à l'API de traitement de langage naturel de Hugging Face pour classer le sentiment, comme le montre la figure~\ref{fig:seq-avis-sentiment}.

\begin{figure}[H]
\centering
\resizebox{\textwidth}{!}{%
\begin{tikzpicture}[>=stealth]
  \colorlet{pboxfill}{amber!10}
  \colorlet{pboxdraw}{amber!80!black}
  \colorlet{actfill}{gray!20}
  \colorlet{actdraw}{gray!60}
  \colorlet{msgcolor}{blue!60!black}
  \tikzset{
    participant/.style={draw=pboxdraw, fill=pboxfill, rounded corners=2pt, minimum width=2.65cm, minimum height=0.65cm, align=center, font=\bfseries\small},
    activation/.style={draw=actdraw, fill=actfill},
    msg/.style={thick, msgcolor},
    ret/.style={thick, dashed, gray!80!black},
    note/.style={midway, above, fill=white, inner sep=1.5pt, font=\scriptsize}
  }

  \foreach \x/\name in {0/Client,3.65/Portail Client,7.3/API Backend,10.95/Celery,14.6/Analyse} {
    \node[participant] at (\x,0) {\name};
    \draw[dashed, thick, gray!55] (\x,-0.45) -- (\x,-8.3);
  }

  \draw[activation] (-0.09,-0.75) rectangle (0.09,-1.2);
  \draw[activation] (3.56,-1.0) rectangle (3.74,-4.0);
  \draw[activation] (7.21,-1.9) rectangle (7.39,-4.45);
  \draw[activation] (10.86,-2.85) rectangle (11.04,-7.35);
  \draw[activation] (14.51,-4.75) rectangle (14.69,-6.0);
  \draw[activation] (7.21,-6.85) rectangle (7.39,-8.05);

  \draw[->, msg] (0.12,-1.0) -- node[note] {1. Avis et note} (3.53,-1.0);
  \draw[->, msg] (3.77,-1.95) -- node[note] {2. Créer avis} (7.18,-1.95);
  \draw[->, msg] (7.42,-2.95) -- node[note] {3. Tâche Celery} (10.83,-2.95);
  \draw[<-, ret] (3.77,-3.85) -- node[note] {4. Avis créé} (7.18,-3.85);
  \draw[->, msg] (11.07,-4.95) -- node[note] {5. Inférence IA} (14.48,-4.95);
  \draw[<-, ret] (11.07,-5.85) -- node[note] {6. Label + score} (14.48,-5.85);
  \draw[->, msg] (10.83,-6.9) -- node[note] {7. Sauvegarder sentiment} (7.42,-6.9);
  \draw[->, msg] (7.42,-7.65) -- ++(0.75,0) -- ++(0,-0.45) -- node[right, align=left, fill=white, inner sep=1.5pt, font=\scriptsize, pos=0.55] {8. Mise à jour\\stats gérant} (7.42,-8.1);
\end{tikzpicture}
}
\caption{Diagramme de séquence de la création d'un avis et analyse de sentiment}
\label{fig:seq-avis-sentiment}
\end{figure}

La figure~\ref{fig:seq-avis-sentiment} montre que la création de l'avis reste rapide pour le client, tandis que l'analyse est exécutée ensuite par Celery puis sauvegardée pour alimenter les statistiques.

\section{Modèle logique de données}
Le modèle logique de données est construit à partir des modèles Django de la codebase. Le tableau~\ref{tab:mld-resume} ne liste pas toutes les colonnes techniques. Il résume les tables structurantes et leurs relations.

\begin{longtable}{p{3.2cm}p{4.2cm}p{6.2cm}}
\caption{Synthèse du modèle logique de données}
\label{tab:mld-resume}\\
\toprule
\textbf{Table / entité} & \textbf{Relations principales} & \textbf{Rôle} \\
\midrule
\endfirsthead
\toprule
\textbf{Table / entité} & \textbf{Relations principales} & \textbf{Rôle} \\
\midrule
\endhead
\code{Utilisateur} & Rôle applicatif. & Compte de connexion et séparation gérant, serveur, cuisinier, client. \\
\code{Categorie}, \code{Plat} & Catégorie 1..N plats. & Structure de la carte, prix, disponibilité, image et temps de préparation. \\
\code{Ingredient}, \code{PlatIngredient} & Relation N..M entre plats et ingrédients. & Recettes et quantités requises pour relier menu et stock. \\
\code{Table}, \code{PlanText} & Table liée aux commandes et réservations. & Plan de salle, capacité, position et statut opérationnel. \\
\code{Commande}, \code{CommandeLigne} & Commande 1..N lignes ; table, serveur ou client optionnels. & Prise de commande, suivi cuisine, prix figé et total. \\
\code{Reservation} & Client, table, date et créneau. & Réservation avec statut et contrôle de disponibilité. \\
\code{Paiement}, \code{PaiementItem} & Paiement lié à une commande ; contributions liées aux lignes. & Encaissement et partage possible de l'addition. \\
\code{Avis}, \code{Analyse}\\\code{Sentiment} & Avis 1..1 analyse ; avis lié à plat ou commande. & Commentaire client, note optionnelle, label et score de sentiment. \\
\code{LoyaltyProfile}, \code{Reward} & Profil 1..1 utilisateur ; transactions et récompenses. & Programme de fidélité. \\
\code{WeatherData} & Données datées. & Support aux indicateurs et prévisions simples. \\
\code{Restaurant}\\\code{Configuration} & Configuration unique. & Paramètres du restaurant, identité, devise et réglages de service. \\
\bottomrule
\end{longtable}

\section{Règles métier principales}
Les règles suivantes structurent le comportement attendu :

\begin{itemize}
  \item un utilisateur possède un rôle qui détermine les routes et les actions accessibles ;
  \item une commande peut être liée à une table, à un serveur et parfois à un client ;
  \item chaque ligne de commande conserve le prix unitaire du plat au moment de la commande ;
  \item les statuts de lignes pilotent le KDS : attente, préparation, prêt, servi ou annulé ;
  \item un paiement est toujours lié à une commande et peut comporter des contributions par ligne ;
  \item un avis appartient à un utilisateur et peut être lié à un plat ou à une commande ;
  \item une analyse de sentiment est associée à un seul avis ;
  \item le stock repose sur les ingrédients, les seuils d'alerte et les mouvements.
\end{itemize}

\section{Architecture fonctionnelle cible}
La conception sépare deux usages : une application staff pour les rôles internes et un portail client. Les deux utilisent le même backend API.

\begin{figure}[H]
\centering
\fbox{\begin{minipage}{0.94\textwidth}
\centering
\schemaNode{Application staff\\Gérant, serveur, cuisinier}
\schemaArrow
\schemaNode{Backend Django REST Framework\\API REST, permissions}
\schemaArrow
\schemaNode{MySQL\\Données métier}

\vspace{0.45cm}

\schemaNode{Portail client\\Menu, réservation, paiement}
\schemaArrow
\schemaNode{Redis\\WebSocket et broker}
\schemaArrow
\schemaNode{Celery + Hugging Face\\Tâches et sentiment}
\end{minipage}}
\caption{Architecture fonctionnelle générale de Tastify}
\label{fig:architecture-fonctionnelle}
\end{figure}

La figure~\ref{fig:architecture-fonctionnelle} résume le choix principal du projet : un backend commun expose les services, tandis que les interfaces restent séparées selon les usages.

\section*{Conclusion}
\addcontentsline{toc}{section}{Conclusion}
L'analyse des besoins fait apparaître un système multi-acteurs, organisé autour d'un backend commun et de deux interfaces web. Les diagrammes existants donnent une base UML, complétée par le modèle logique tiré du code. Le chapitre suivant détaille l'analyse de sentiment, utilisée pour traiter les retours clients.
